\section{İlgili çalışmalar}\label{sec:related}

\paragraph{Min--maks mTSP} Çoklu gezgin satıcı problemi ve
formülasyonları \citet{bektas2006} tarafından derlenmiştir. Min--maks (tamamlanma süresi) varyantı,
amaç fonksiyonu düz olduğundan yerel arama açısından min-toplam varyantından daha zordur:
hamlelerin çoğu en uzun rotayı değiştirmediği için amaç değerini de
etkilemez. \citet{necula2015} tarafından önerilen 16 örnekli
mTSPLib kümesi (eil51, berlin52, eil76, rat99;
$k\in\{2,3,5,7\}$) standart küçük kıyaslama kümesi hâline gelmiştir. Son dönem literatüründe öne çıkan
metasezgiseller memetik ve yinelemeli yerel arama algoritmalarıdır: He ve Hao'nun
komşuluk indirgemeli hibrit araması (HSNR) \citep{hehao2022} ile kenar birleştirmeli
çaprazlama kullanan memetik algoritması (MA) \citep{hehao2023},
\citet{mahmoudinazlou2024} tarafından önerilen hibrit genetik algoritma (HGA),
\citet{hehaoxia2024} tarafından önerilen öğrenme güdümlü
yinelemeli yerel arama (MILS), yinelemeli
iki aşamalı sezgisel algoritma (ITSHA) \citep{zheng2022} ve en yeni örnek olarak
pekiştirmeli öğrenmeyle desteklenen bir kur--birleştir--çöz--uyarla (CMSA) şeması \citep{rodriguez2026}.
Bu yöntemler, uzun bütçeler altında çözüm kalitesini önceleyecek biçimde tasarlanmıştır:
\citet{hehaoxia2024} çalışmasındaki büyük örnek kümesi $\mathbb{L}$ ($n=1379$--$5915$),
koşu başına $(n/100)\times 4$ dakikalık, yani en büyük örnek için yaklaşık dört
saatlik bir kesme süresiyle ve örnek başına 20 koşu yapılarak çözülmektedir.

\paragraph{LKH-3} LKH-3 \citep{helsgaun2017}, gezgin satıcı problemi (TSP) için geliştirilmiş
Lin--Kernighan--Helsgaun çözücüsünü \citep{helsgaun2000}, bir mTSP'yi $k$ depo kopyası içeren
bir TSP'ye dönüştürerek ve amacı bir ceza fonksiyonuyla ele alarak kısıtlı rotalama
problemlerine genişletir; MINMAX, desteklediği mTSP amaçlarından biridir. Olgunluğu,
herkese açık olması ve güçlü sonuçları, LKH-3'ü rotalama varyantları için yaygın olarak kullanılan
genel amaçlı bir referans ve uygulayıcılar için doğal bir karşılaştırma tabanı
hâline getirmiştir. Ancak bu çalışmada raporlanan deneylerde bu çözücünün min--maks kipi, binlerce hedef içeren
örneklerde operasyonel zaman bütçeleri içinde çoğu zaman rekabetçi bir çözüm
döndürememiştir (bkz.\ Kıyaslama sonuçları bölümü).

\paragraph{Önce rota, sonra küme (route-first/cluster-second)} Tek bir dev turun araç
rotalarına bölünmesi fikri \citet{beasley1983} çalışmasına dayanır; \citet{prins2004},
optimal bölmeyi araç rotalama problemi (VRP) için geliştirdiği bir genetik algoritmanın çekirdeği hâline getirmiştir. Min--maks
$k$-TSP için tur bölme, \citet{frederickson1978} tarafından verilen
$(5/2-1/k)$ oranlı klasik yaklaşıklık algoritmasının da temelini oluşturur.
PSB, Greedy-Edge dev turu \citep{bentley1992} üzerinde optimal min--maks bölme
(Lemma~\ref{lem:split}) uygulamakta ve bunu rotalar arası dengeleme ile
\citet{shaw1998} yaklaşımından esinlenen bir geniş komşuluk aramasıyla birleştirmektedir.

\paragraph{Hasat robotiği ve tarla lojistiği} Hasat, ilaçlama ve tarla gözlemi için kullanılan tarla robotları
\citet{bechar2016}, tarımsal drone'lar ise \citet{rejeb_2022_drones} tarafından
derlenmiştir. Robotik elma
hasat makineleri ticari bahçelerde, meyve başına yaklaşık 6--7\,s çevrim süresi
ve 65--84\% koparma başarısıyla değerlendirilmiştir \citep{silwal_2017_design,
hohimer_2019_design,zhang_2024_automated,zhu_2025_advancement}; derin öğrenme tabanlı
dedektörler elmaları 0.84--0.93 kesinlik (precision) ve
0.82--0.97 yakalama oranı (recall) ile konumlandırmaktadır \citep{kang_2020_fast,sapkota_2024_comparing}; havadan kavrama ve toplama
konusundaki çalışmalar \citet{ruggiero_2018_aerial} tarafından derlenmiş, uygulaması ise
\citet{kumar_2024_design} tarafından gösterilmiştir. Hasat destek platformları yüksek yoğunluklu bahçelerde
toplama verimliliğini artırmaktadır \citep{baugher_2009_mobile,zhang_2016_design}. Tarla lojistiğinde
araç rotalama problemi \citet{bochtis2009} tarafından formüle edilmiştir;
bahçelerde ise hedefler kusurlu bir kafes üzerindeki ağaçlardır, robotlar baş tarladaki bir
üs noktasından yola çıkar ve araç üstü algılama, olgun meyveye ilişkin kaba ön kestirimlerin yerini aldıkça
hedef kümesi operasyon boyunca değişir. Bu nedenle planlama hızı iki kez
önem kazanır: her planda bir kez, her \emph{yeniden} planlamada bir kez daha. Hakem değerlendirmesindeki tamamlayıcı çalışmamızda bu tür kafes benzeri nokta kümeleri için hızlı
yapıcı sezgiseller ve eksene duyarlı paralel bölümleme analiz edilmiştir \citep{asil2026frame};
bu makale, söz konusu çalışmanın Greedy-Edge'in en dayanıklı
yapıcı yöntem olduğu bulgusunu temel almakta ve o çalışmada ele alınmayan depolu, çok robotlu
min--maks amacını incelemektedir.
